\section*{Problems}

\subsection{Problem Statements}

\subsubsection*{1. Fundamental Level (Simple / Direct)}

\begin{problem}
	\label{prob:2.5.1}
	\sidenote{\href{https://stempunkxyz.wordpress.com/2021/08/02/teorema-de-bayes-1/}{Solution to Problem \ref{prob:2.5.1}}}
	Find the probability that a single roll of a fair 6-sided die results in a number less than 4 if:
	\begin{enumerate}
		\item no further information is available;
		\item it is known with certainty that the roll resulted in an odd number.
	\end{enumerate}
\end{problem}

\begin{problem}
	\label{prob:2.5.6}
	In a group of statistics students, 60\% are female ($M$) and 40\% are male ($H$). It is known that 10\% of the females are left-handed ($Z$) and 15\% of the males are left-handed. If a person is selected at random and turns out to be left-handed, use Bayes' Theorem to calculate the probability that the person is female.
\end{problem}

\begin{problem}
	\label{prob:2.5.4}
	A student is preparing for a comprehensive examination. If the student studies properly ($E$), the probability of passing ($A$) is 0.90. If the student does not study ($E'$), the probability of passing is only 0.30. The probability that the student studies is 0.70.
	\begin{enumerate}
		\item What is the total probability that the student passes the exam?
		\item If the student passed, what is the conditional probability that they studied?
		\item If the student did not pass, what is the probability that they did not study?
	\end{enumerate}
\end{problem}

\subsubsection*{2. Operational Level (Intermediate / Developmental)}

\begin{problem}
	\label{prob:2.5.7}
	Box 1 contains 3 red marbles and 2 blue marbles, while Box 2 contains 2 red marbles and 8 blue marbles.
	Two fair dice are rolled and the sum of their faces is calculated: if a sum of at most six points is obtained, a marble is chosen from Box 1; otherwise, a marble is chosen from Box 2.
	\begin{enumerate}
		\item Calculate the marginal probability that the chosen marble is red.
		\item If the person who rolled the dice does not reveal the sum obtained, but does show that the chosen marble was red, what is the posterior probability that it was drawn from Box 1?
	\end{enumerate}
\end{problem}

\begin{problem}
	\label{prob:2.5.2}
	A clinical diagnostic medical test to detect a pathology has the following conditional metrics:
	\begin{itemize}
		\item Sensitivity ($P(+|E)$): If the person has the disease, the test is positive with probability 0.95.
		\item Specificity ($P(-|E')$): If the person does not have the disease, the test is negative with probability 0.90.
	\end{itemize}
	If the prevalence of the disease in the population is 2\% ($P(E)=0.02$), calculate:
	\begin{enumerate}
		\item The probability that a person selected at random obtains a positive result on the test.
		\item The Positive Predictive Value (PPV): the probability that a person who tested positive actually has the disease.
		\item The Negative Predictive Value (NPV): the probability that a person who tested negative is actually healthy.
	\end{enumerate}
\end{problem}

\begin{problem}
	\label{prob:2.5.5}
	A digital communication channel transmits binary signals (0 or 1). Due to noise in the channel, the probability that a transmitted 0 bit is erroneously received as a 1 is 0.02, and the probability that a transmitted 1 bit is received as a 0 is 0.03. If 60\% of the bits transmitted by the source are 0:
	\begin{enumerate}
		\item What is the total probability that a received bit is 1?
		\item If a 1 is registered at the receiver, what is the exact probability that a 1 was actually transmitted by the source?
	\end{enumerate}
\end{problem}

\subsubsection*{3. Analytical Level (Advanced / Connection)}

\begin{problem}
	\label{prob:2.5.3}
	Three assembly lines $M_1, M_2, M_3$ account for 50\%, 30\%, and 20\% of the total semiconductor production of an industrial plant, respectively. The historical rates of defective parts are 3\%, 5\%, and 7\% for each assembly line.
	\begin{enumerate}
		\item If the auditing department selects a part at random and it is defective, calculate the complete vector of posterior probabilities to determine from which line the part originated.
		\item Which of the three lines is statistically the most probable cause of the defective part? Formally justify your answer.
	\end{enumerate}
\end{problem}

\begin{problem}
	\label{prob:2.5.8}
	Rigorously prove Bayes' Theorem for a finite and exhaustive partition $\{B_1, B_2, \dots, B_k\}$ of the sample space $S$ given an observable event $A$ with $P(A)>0$:
	\begin{align}
		P(B_j | A) = \frac{P(A|B_j)P(B_j)}{\sum_{i=1}^k P(A|B_i)P(B_i)}, \quad \text{for all } j \in \{1, \dots, k\}.
	\end{align}
\end{problem}

\subsubsection*{4. Challenging Level (Expert / Deep Deduction)}

\begin{problem}
	\label{prob:2.5.9}
	\textbf{The Prosecutor's Fallacy and Odds Ratio Formulation:} In a forensic investigation, a biological sample is recovered from a crime scene. The extracted DNA profile has a random match probability in the general population of $\alpha = 10^{-6}$ (i.e., $P(\text{match} | \text{innocent}) = 10^{-6}$). A suspect is identified through a mass database search and their DNA perfectly matches the sample ($P(\text{match} | \text{guilty}) = 1$). The prosecutor claims to the jury that, since the random match probability is barely one in a million, the probability that the suspect is innocent is only $0.0001\%$.
	
	Using the odds ratio formulation of Bayes' Theorem, prove that the prosecutor's reasoning is flawed. If $N = 5,000,000$ people with potential access to the scene live in the region and no other evidence against the suspect exists, calculate the true posterior probability of guilt $P(\text{guilty} | \text{match})$.
\end{problem}

\begin{problem}
	\label{prob:2.5.10}
	\textbf{Bayesian Spam Filter (Naive Bayes Classifier):} A mail server processes an incoming stream where 60\% of messages are unsolicited spam ($S$) and 40\% are legitimate email ($S'$). An algorithm analyzes the presence of three keywords ($W_1, W_2, W_3$). Under the conditional independence assumption (Naive Bayes), the appearance of the words given the email class is independent:
	\begin{align}
		P(W_1 \cap W_2 \cap W_3 | S) &= P(W_1|S) \cdot P(W_2|S) \cdot P(W_3|S) \\
		P(W_1 \cap W_2 \cap W_3 | S') &= P(W_1|S') \cdot P(W_2|S') \cdot P(W_3|S')
	\end{align}
	The empirical conditional probabilities in spam are: $P(W_1|S)=0.80$, $P(W_2|S)=0.70$, $P(W_3|S)=0.60$; and in legitimate email they are: $P(W_1|S')=0.10$, $P(W_2|S')=0.20$, $P(W_3|S')=0.10$.
	
	Prove the general decision rule of the Naive Bayes classifier and calculate the exact posterior probability that an email simultaneously containing all three words ($W = W_1 \cap W_2 \cap W_3$) is actually spam.
\end{problem}


\subsection{Detailed Solutions}

\subsubsection*{1. Fundamental Level (Simple / Direct)}

\begin{solproblem}
	\begin{enumerate}
		\item \textbf{Without information:} The complete sample space has 6 equally likely outcomes $S=\{1,2,3,4,5,6\}$.
		The outcomes strictly less than 4 form the subset $A=\{1,2,3\}$.
		Therefore, $P(A) = \frac{\card{A}}{\card{S}} = \frac{3}{6} = \frac{1}{2} = 0.50$.
		
		\item \textbf{With odd number information:} The new conditioned sample space is the event $I=\{1,3,5\}$ with $\card{I}=3$.
		The elements of $A$ that are also in $I$ form the intersection $A \cap I = \{1,3\}$, with $\card{A \cap I}=2$.
		Applying the definition of conditional probability (or Bayes simplified by equal probability):
		\begin{align}
			P(A|I) = \frac{P(A \cap I)}{P(I)} = \frac{2/6}{3/6} = \frac{2}{3} \approx 0.6667.
		\end{align}
	\end{enumerate}
\end{solproblem}

\begin{solproblem}
	Let the events be $M$ (female) with $P(M)=0.60$ and $H$ (male) with $P(H)=0.40$.
	The left-handed rates are $P(Z|M)=0.10$ and $P(Z|H)=0.15$.
	
	By the law of total probability, the probability that a person chosen at random is left-handed is:
	\begin{align}
		P(Z) &= P(Z|M)P(M) + P(Z|H)P(H) = (0.10)(0.60) + (0.15)(0.40) \\
		&= 0.060 + 0.060 = 0.120 \text{ (or 12\%).}
	\end{align}
	By Bayes' Theorem, the posterior probability that a left-handed person is female is:
	\begin{align}
		P(M|Z) = \frac{P(Z|M)P(M)}{P(Z)} = \frac{0.060}{0.120} = \frac{1}{2} = 0.50 \text{ (50\%).}
	\end{align}
\end{solproblem}

\begin{solproblem}
	Let $A$ be the event of passing and $E$ the event of studying ($P(E)=0.70$, $P(E')=0.30$).
	The conditional probabilities of passing are $P(A|E)=0.90$ and $P(A|E')=0.30$.
	\begin{enumerate}
		\item Total marginal probability of passing:
		\begin{align}
			P(A) &= P(A|E)P(E) + P(A|E')P(E') = (0.90)(0.70) + (0.30)(0.30) \\
			&= 0.63 + 0.09 = 0.72 \text{ (72\%).}
		\end{align}
		
		\item By Bayes' Theorem, if the student passed, the probability that they studied is:
		\begin{align}
			P(E|A) = \frac{P(A|E)P(E)}{P(A)} = \frac{0.63}{0.72} = \frac{7}{8} = 0.875 \text{ (87.5\%).}
		\end{align}
		
		\item For the event of failing ($A'$), we have $P(A') = 1 - 0.72 = 0.28$.
		The conditional probabilities of failing given whether they studied or not are $P(A'|E)=1-0.90=0.10$ and $P(A'|E')=1-0.30=0.70$.
		By Bayes' Theorem applied to the complements:
		\begin{align}
			P(E'|A') = \frac{P(A'|E')P(E')}{P(A')} = \frac{(0.70)(0.30)}{0.28} = \frac{0.21}{0.28} = \frac{3}{4} = 0.75 \text{ (75\%).}
		\end{align}
	\end{enumerate}
\end{solproblem}

\subsubsection*{2. Operational Level (Intermediate / Developmental)}

\begin{solproblem}
	The sample space for rolling two dice has 36 points. A sum of at most 6 includes outcomes with sums 2, 3, 4, 5, and 6, which counts $1+2+3+4+5 = 15$ simple events.
	Thus, the probability of selecting Box 1 is $P(C_1) = \frac{15}{36} = \frac{5}{12}$, and selecting Box 2 is $P(C_2) = 1 - \frac{5}{12} = \frac{7}{12}$.
	The proportions of red marbles ($R$) in the boxes are $P(R|C_1) = \frac{3}{5}$ and $P(R|C_2) = \frac{2}{10} = \frac{1}{5}$.
	
	\begin{enumerate}
		\item Total probability of drawing a red marble:
		\begin{align}
			P(R) &= P(R|C_1)P(C_1) + P(R|C_2)P(C_2) = \left(\frac{3}{5} \cdot \frac{5}{12}\right) + \left(\frac{1}{5} \cdot \frac{7}{12}\right) \\
			&= \frac{15}{60} + \frac{7}{60} = \frac{22}{60} = \frac{11}{30} \approx 0.3667.
		\end{align}
		
		\item Inverse conditional probability (Bayes) that it came from Box 1:
		\begin{align}
			P(C_1|R) = \frac{P(R|C_1)P(C_1)}{P(R)} = \frac{15/60}{22/60} = \frac{15}{22} \approx 0.6818 \text{ (68.18\%).}
		\end{align}
	\end{enumerate}
\end{solproblem}

\begin{solproblem}
	With prevalence $P(E)=0.02$ (and $P(E')=0.98$), sensitivity $P(+|E)=0.95$, and specificity $P(-|E')=0.90$ (hence the false positive rate is $P(+|E')=1-0.90=0.10$):
	\begin{enumerate}
		\item Total probability of a positive test result:
		\begin{align}
			P(+) &= P(+|E)P(E) + P(+|E')P(E') = (0.95)(0.02) + (0.10)(0.98) \\
			&= 0.019 + 0.098 = 0.117 \text{ (11.7\%).}
		\end{align}
		
		\item Positive Predictive Value (PPV) by Bayes' Theorem:
		\begin{align}
			P(E|+) = \frac{P(+|E)P(E)}{P(+)} = \frac{0.019}{0.117} = \frac{19}{117} \approx 0.1624 \text{ (16.24\%).}
		\end{align}
		\emph{Interpretation:} Despite the test's high sensitivity (95\%), due to low disease prevalence, 83.76\% of positive results are false positives.
		
		\item Negative Predictive Value (NPV):
		\begin{align}
			P(-) &= 1 - P(+) = 1 - 0.117 = 0.883 \\
			P(E'|-) &= \frac{P(-|E')P(E')}{P(-)} = \frac{(0.90)(0.98)}{0.883} = \frac{0.882}{0.883} \approx 0.9989 \text{ (99.89\%).}
		\end{align}
	\end{enumerate}
\end{solproblem}

\begin{solproblem}
	Let $T_0, T_1$ denote transmission and $R_0, R_1$ reception, with $P(T_0)=0.60$ and $P(T_1)=0.40$.
	Channel probabilities are $P(R_1|T_0)=0.02$ (error) and $P(R_1|T_1)=1-0.03=0.97$ (success).
	\begin{enumerate}
		\item Marginal probability of receiving bit 1:
		\begin{align}
			P(R_1) &= P(R_1|T_0)P(T_0) + P(R_1|T_1)P(T_1) = (0.02)(0.60) + (0.97)(0.40) \\
			&= 0.012 + 0.388 = 0.400 \text{ (40\%).}
		\end{align}
		
		\item Posterior probability that a received 1 is authentic:
		\begin{align}
			P(T_1|R_1) = \frac{P(R_1|T_1)P(T_1)}{P(R_1)} = \frac{0.388}{0.400} = 0.970 \text{ (97\%).}
		\end{align}
	\end{enumerate}
\end{solproblem}

\subsubsection*{3. Analytical Level (Advanced / Connection)}

\begin{solproblem}
	Let $D$ be a defective part and the three production lines be $M_1, M_2, M_3$.
	\begin{enumerate}
		\item First, we find the marginal probability of a defect across the entire plant:
		\begin{align}
			P(D) &= P(D|M_1)P(M_1) + P(D|M_2)P(M_2) + P(D|M_3)P(M_3) \\
			&= (0.03)(0.50) + (0.05)(0.30) + (0.07)(0.20) = 0.015 + 0.015 + 0.014 = 0.044.
		\end{align}
		We calculate the three posterior probabilities (generalized Bayes' Theorem):
		\begin{align}
			P(M_1|D) &= \frac{0.015}{0.044} = \frac{15}{44} \approx 0.3409 \text{ (34.09\%)} \\
			P(M_2|D) &= \frac{0.015}{0.044} = \frac{15}{44} \approx 0.3409 \text{ (34.09\%)} \\
			P(M_3|D) &= \frac{0.014}{0.044} = \frac{14}{44} = \frac{7}{22} \approx 0.3182 \text{ (31.82\%)}
		\end{align}
		
		\item Lines $M_1$ and $M_2$ share exactly the highest conditional probability at 34.09\% each. Although $M_1$ has a very low defect rate (3\%) compared to $M_2$ (5\%) or $M_3$ (7\%), its massive production volume (50\% of the plant's total) means it matches line $M_2$ in the absolute number of defective parts generated.
	\end{enumerate}
\end{solproblem}

\begin{solproblem}
	Let $\{B_1, B_2, \dots, B_k\}$ be an exhaustive partition ($\bigcup_{i=1}^k B_i = S$) that is mutually exclusive ($B_i \cap B_m = \emptyset$ for all $i \neq m$).
	By the multiplication rule of conditional probability, for any particular index $j$:
	\begin{align}
		P(B_j \cap A) = P(B_j | A) P(A) = P(A | B_j) P(B_j).
	\end{align}
	Solving for the inverse conditional probability $P(B_j|A)$ given that $P(A)>0$:
	\begin{align}
		P(B_j | A) = \frac{P(A | B_j) P(B_j)}{P(A)}. \label{eq:bayes_num}
	\end{align}
	On the other hand, since the events $B_i$ form a partition of $S$, we can express event $A$ as the disjoint union of its intersections with each cell of the partition:
	\begin{align}
		A = A \cap S = A \cap \left( \bigcup_{i=1}^k B_i \right) = \bigcup_{i=1}^k (A \cap B_i).
	\end{align}
	Applying Kolmogorov's third axiom (additivity for disjoint events), the total probability of $A$ is:
	\begin{align}
		P(A) = \sum_{i=1}^k P(A \cap B_i) = \sum_{i=1}^k P(A | B_i) P(B_i). \label{eq:prob_total_k}
	\end{align}
	Substituting the denominator (\ref{eq:prob_total_k}) into equation (\ref{eq:bayes_num}) rigorously yields the general formula of Bayes' Theorem for $k$ states.
\end{solproblem}

\subsubsection*{4. Challenging Level (Expert / Deep Deduction)}

\begin{solproblem}
	Let $I$ be the event of being the true perpetrator of the crime, and $E$ the event of a DNA profile match.
	The prosecutor commits a conditional probability fallacy by erroneously equating $P(E|I') = 10^{-6}$ with the inverse probability $P(I'|E)$.
	
	Let us evaluate formal validity using the odds ratio version of Bayes' Theorem (Odds Form):
	\begin{align}
		\frac{P(I|E)}{P(I'|E)} = \underbrace{\frac{P(E|I)}{P(E|I')}}_{\text{Likelihood Ratio (LR)}} \times \underbrace{\frac{P(I)}{P(I')}}_{\text{Prior Odds}}.
	\end{align}
	With no other evidence besides geographic presence in a population of $N = 5,000,000$ inhabitants, the prior probability that an arbitrary individual is guilty is $P(I) = \frac{1}{N} = \frac{1}{5,000,000} = 2 \times 10^{-7}$, while $P(I') \approx 1$.
	
	The likelihood ratio of the evidence is:
	\begin{align}
		\text{LR} = \frac{P(E|I)}{P(E|I')} = \frac{1}{10^{-6}} = 1,000,000.
	\end{align}
	We now calculate the posterior odds by multiplying the likelihood ratio by the prior odds:
	\begin{align}
		\text{Odds}(I|E) = 1,000,000 \times \frac{1 / 5,000,000}{1} = \frac{1,000,000}{5,000,000} = \frac{1}{5} = 0.20.
	\end{align}
	Converting the posterior odds $\frac{P(I|E)}{1-P(I|E)} = \frac{1}{5}$ to exact probability:
	\begin{align}
		P(I|E) = \frac{0.20}{1 + 0.20} = \frac{1/5}{6/5} = \frac{1}{6} \approx 0.1667 \text{ (16.67\%).}
	\end{align}
	\textbf{Proof and refutation:} The probability that the suspect is guilty is only 16.67\% (and not 99.9999\% as claimed by the prosecutor). The flaw lies in ignoring the fact that testing a profile across a population of 5 million where the random match frequency is 1 in 1 million statistically yields about 5 purely coincidental false positives in the population.
\end{solproblem}

\begin{solproblem}
	Let the observed keyword vector be $W = W_1 \cap W_2 \cap W_3$. By Bayes' Theorem for dichotomous classes ($S$ and $S'$):
	\begin{align}
		P(S|W) = \frac{P(W|S)P(S)}{P(W|S)P(S) + P(W|S')P(S')}.
	\end{align}
	Under the conditional independence hypothesis (Naive Bayes), we replace the joint probability with the product of conditional marginals:
	\begin{align}
		P(W|S) &= P(W_1|S) \cdot P(W_2|S) \cdot P(W_3|S) = (0.80)(0.70)(0.60) = 0.336. \\
		P(W|S') &= P(W_1|S') \cdot P(W_2|S') \cdot P(W_3|S') = (0.10)(0.20)(0.10) = 0.002.
	\end{align}
	With the server's prior proportions $P(S)=0.60$ and $P(S')=0.40$, we weight both terms:
	\begin{align}
		P(W|S)P(S) &= (0.336)(0.60) = 0.2016 \\
		P(W|S')P(S') &= (0.002)(0.40) = 0.0008
	\end{align}
	The total marginal probability of receiving this exact set of words in any email is:
	\begin{align}
		P(W) = 0.2016 + 0.0008 = 0.2024.
	\end{align}
	Therefore, the posterior probability that the incoming message is spam is:
	\begin{align}
		P(S|W) = \frac{0.2016}{0.2024} = \frac{252}{253} \approx 0.996047 \text{ (99.60\%).}
	\end{align}
	The Naive Bayes classifier successfully filters and categorizes the message as spam with over 99.6\% certainty.
\end{solproblem}
